\documentclass[10pt, letterpaper, twocolumn]{article}

\usepackage[utf8]{inputenc}
\usepackage[T1]{fontenc}
\usepackage{times}
\usepackage{amsmath, amssymb}
\usepackage{textcomp}
\usepackage{microtype}
\usepackage[top=1in, bottom=1in, left=0.75in, right=0.75in]{geometry}
\usepackage{hyperref}
\usepackage{natbib}
\usepackage{url}
\usepackage{booktabs}
\usepackage{longtable}
\usepackage{array}
\usepackage{calc}
\makeatletter
\providecommand{\real}[1]{\dimexpr#1\p@\relax}
\makeatother

% Unicode character support
\DeclareUnicodeCharacter{00D7}{\texttimes}
\DeclareUnicodeCharacter{2264}{\ensuremath{\leq}}
\DeclareUnicodeCharacter{2265}{\ensuremath{\geq}}
\DeclareUnicodeCharacter{2248}{\ensuremath{\approx}}
\usepackage{graphicx}
\usepackage{titlesec}
\usepackage{abstract}
\usepackage{fancyhdr}
\usepackage{xcolor}

% Hyperlink styling
\hypersetup{
    colorlinks=true,
    linkcolor=black,
    citecolor=black,
    urlcolor={blue!70!black}
}

% Section formatting - bold, compact
\titleformat{\section}{\large\bfseries}{\thesection}{1em}{}
\titleformat{\subsection}{\normalsize\bfseries}{\thesubsection}{1em}{}
\titleformat{\subsubsection}{\normalsize\itshape}{\thesubsubsection}{1em}{}
\titlespacing*{\section}{0pt}{2.5ex plus 1ex minus .2ex}{1ex plus .2ex}
\titlespacing*{\subsection}{0pt}{2ex plus 1ex minus .2ex}{0.75ex plus .2ex}

% Abstract styling
\renewcommand{\abstractnamefont}{\normalfont\bfseries}
\renewcommand{\abstracttextfont}{\normalfont\small}
\setlength{\absleftindent}{0.5in}
\setlength{\absrightindent}{0.5in}

% Header/footer
\pagestyle{fancy}
\fancyhf{}
\renewcommand{\headrulewidth}{0pt}
\fancyhead[C]{\small\textit{Proceedings of the Institute for a Christian Machine Intelligence}}
\fancyfoot[C]{\small\thepage}

% First page style
\fancypagestyle{firstpage}{
  \fancyhf{}
  \renewcommand{\headrulewidth}{0pt}
  \fancyfoot[C]{\small\thepage}
}

% Compact lists
\usepackage{enumitem}
\setlist{nosep, leftmargin=*}

% Pandoc tightlist support
\providecommand{\tightlist}{\setlength{\itemsep}{0pt}\setlength{\parskip}{0pt}}

% Paragraph spacing
\setlength{\parindent}{1em}
\setlength{\parskip}{0pt}

% Title block
\title{\Large\textbf{biblical-render: Design and Validation of a Biblical Text Style Transfer Tool}\\[0.75em]
\normalsize\textit{ICMI Working Paper D}}
\author{Institute for a Christian Machine Intelligence}
\date{March 2026}

\begin{document}

\twocolumn[
  \begin{@twocolumnfalse}
    \maketitle
    \thispagestyle{firstpage}
        \begin{abstract}
    This document describes the design, implementation, and validation of `biblical-render`, a CLI tool that transforms arbitrary modern prose into Biblical scripture format across 15 translation styles --- 8 English Bible translations and 7 historical languages. The tool is intended as a research instrument for studying whether the formal linguistic and structural properties of Biblical text exert an independent aligning effect on language model behavior. We present the system architecture, prompting methodology, and a systematic validation of output quality across all supported styles using a common reference text.
    \end{abstract}
        \vspace{1.5em}
  \end{@twocolumnfalse}
]

\section*{1. Introduction}\label{introduction}

\subsection*{1.1 Research Context}\label{research-context}

Large language models are known to be sensitive to the stylistic and structural properties of their input text. Prompt framing, register, and formatting can measurably alter model outputs in ways that extend beyond semantic content. Biblical text represents a distinctive combination of formal properties: archaic or elevated register, paratactic sentence structure, chapter-verse segmentation, parallelism, and a tone of moral authority. Whether these formal properties alone --- independent of theological content --- have an aligning effect on model behavior is an open empirical question.

\subsection*{1.2 Purpose of the Tool}\label{purpose-of-the-tool}

\texttt{biblical-render} provides a controlled mechanism for producing Biblical-style renderings of arbitrary secular text. By supporting multiple translation styles, the tool enables researchers to vary the specific linguistic properties of the Biblical framing (e.g., archaic vs.~modern English, formal vs.~colloquial register) while holding semantic content constant. This supports experimental designs that isolate the contribution of individual stylistic features.

\section*{2. System Design}\label{system-design}

\subsection*{2.1 Architecture}\label{architecture}

The tool is a Node.js CLI application that sends text to the Anthropic Messages API (Claude Opus 4.6) with translation-specific system prompts. The architecture is straightforward:

\begin{enumerate}
\def\labelenumi{\arabic{enumi}.}
\tightlist
\item
  \textbf{Input} is accepted as inline text, file path, or stdin pipe
\item
  \textbf{A system prompt} is constructed from a translation-specific style definition plus universal formatting and completeness requirements
\item
  \textbf{The API call} is made with streaming enabled by default
\item
  \textbf{Output} is rendered to stdout and optionally written to a file
\end{enumerate}

\subsection*{2.2 Prompting Strategy}\label{prompting-strategy}

Each translation style is defined by two prompt components:

\textbf{Style definition}: A description of the target translation's specific linguistic features --- pronoun forms, verb constructions, sentence structure, vocabulary sources, tone, and register. These were authored based on the documented translation philosophies and observable characteristics of each Bible translation.

\textbf{Universal requirements}: Applied to all styles, these enforce: - Chapter-verse structural formatting - Completeness (no summarization or abbreviation) - Semantic preservation (no editorializing) - Length matching (output should feel as complete as the input, with Biblical elaboration filling out the text naturally)

For historical language outputs, the prompt additionally requires output in the original script with a literal English gloss after each verse.

\subsection*{2.3 Supported Styles}\label{supported-styles}

The 15 supported styles were selected to span the major axes of variation in Biblical translation:

{\def\LTcaptype{none} % do not increment counter
\begin{tabular}[]{@{}ll@{}}
\toprule\noalign{}
Axis & Styles \\
\midrule\noalign{}
\bottomrule\noalign{}
Formal equivalence (literal) & KJV, ESV, NASB, NKJV \\
Dynamic equivalence (thought-for-thought) & NIV, NLT \\
Paraphrase & MSG \\
Ecclesiastical register & VULGATE \\
Semitic languages & HEBREW, ARAMAIC \\
Koine Greek & GREEK \\
Latin & LATIN \\
African Christian traditions & GEEZ, COPTIC \\
Early Germanic & GOTHIC \\
\end{tabular}
}

This selection covers the formal-dynamic equivalence spectrum in English translations, plus the major historical languages of Biblical transmission.

\section*{3. Validation Methodology}\label{validation-methodology}

\subsection*{3.1 Reference Text}\label{reference-text}

All 15 styles were run against a single reference input: a 274-word prose summary of Anthropic's Claude Constitution, covering its hierarchy of values (safety, ethics, compliance, helpfulness), honesty principles, harm avoidance framework, user wellbeing considerations, and preference for cultivating judgment over rigid rules.

This text was chosen because it is: - Secular and modern in register - Structurally organized (multiple thematic sections) - Conceptually dense but not technical - Relevant to the downstream research application (AI alignment)

\subsection*{3.2 Evaluation Criteria}\label{evaluation-criteria}

Each output was assessed on:

\begin{enumerate}
\def\labelenumi{\arabic{enumi}.}
\tightlist
\item
  \textbf{Structural fidelity}: Does the output use proper chapter-verse formatting?
\item
  \textbf{Stylistic accuracy}: Does the output reflect the distinctive linguistic characteristics of the target translation?
\item
  \textbf{Semantic completeness}: Are all ideas from the input represented without omission or addition?
\item
  \textbf{Register consistency}: Does the output maintain a consistent register throughout?
\item
  \textbf{Script accuracy} (historical languages): Is the output in the correct script with appropriate diacritical marks?
\end{enumerate}

\section*{4. Validation Results}\label{validation-results}

\subsection*{4.1 English Translation Styles}\label{english-translation-styles}

\subsubsection*{KJV (King James Version)}\label{kjv-king-james-version}

\begin{itemize}
\tightlist
\item
  \textbf{Word count}: 696 (2.54x input)
\item
  \textbf{Structural fidelity}: 5 chapters, 22 verses. Correct chapter-verse numbering throughout.
\item
  \textbf{Stylistic markers observed}: Archaic pronouns (``it shall'', ``doth''), inverted syntax (``And there was set forth''), conjunctive openings (``And'', ``For'', ``Behold''), doublets (``nor weaken, nor subvert''), parallelism. The characteristic KJV cadence is sustained throughout.
\item
  \textbf{Semantic completeness}: All five thematic sections of the input are represented. The 1,000-users thought experiment, the anti-paternalism principle, and the judgment-over-rules philosophy are all faithfully rendered.
\item
  \textbf{Assessment}: Strong. The output reads plausibly as KJV scripture.
\end{itemize}

\subsubsection*{NIV (New International Version)}\label{niv-new-international-version}

\begin{itemize}
\tightlist
\item
  \textbf{Word count}: 698 (2.55x input)
\item
  \textbf{Structural fidelity}: 5 chapters, 24 verses.
\item
  \textbf{Stylistic markers observed}: Modern English, clear direct sentences, dignified but accessible. No archaic pronouns. Uses ``For'' and ``Now'' as section transitions rather than ``And it came to pass.''
\item
  \textbf{Semantic completeness}: Full coverage. Notably preserves the nuance of ``almost no exception'' for the lying prohibition.
\item
  \textbf{Assessment}: Strong. Captures the NIV's balance of clarity and gravity.
\end{itemize}

\subsubsection*{ESV (English Standard Version)}\label{esv-english-standard-version}

\begin{itemize}
\tightlist
\item
  \textbf{Word count}: 673 (2.46x input)
\item
  \textbf{Stylistic markers observed}: ``And behold'' retained, formal but readable, chapter subheadings (``The Ordering of Values'', ``The Weighing of Harms''), literary precision. More elevated than NIV but without KJV archaisms.
\item
  \textbf{Semantic completeness}: Full coverage including the judgment-vs-rules distinction.
\item
  \textbf{Assessment}: Strong. The slightly formal literary register is well-differentiated from both KJV and NIV.
\end{itemize}

\subsubsection*{NASB (New American Standard Bible)}\label{nasb-new-american-standard-bible}

\begin{itemize}
\tightlist
\item
  \textbf{Word count}: 539 (1.97x input)
\item
  \textbf{Stylistic markers observed}: Bracketed supplied words (``{[}one would inscribe{]}'', ``{[}the virtues{]}''), stiff formal constructions mirroring source language structure, precise vocabulary. Notably more compressed than other styles, consistent with NASB's literal philosophy.
\item
  \textbf{Semantic completeness}: Full coverage, though more terse in elaboration.
\item
  \textbf{Assessment}: Strong. The bracketed interpolations and formal stiffness are distinctively NASB.
\end{itemize}

\subsubsection*{MSG (The Message)}\label{msg-the-message}

\begin{itemize}
\tightlist
\item
  \textbf{Word count}: 682 (2.49x input)
\item
  \textbf{Stylistic markers observed}: Conversational tone (``listen carefully here''), contractions (``don't'', ``here's''), colloquial asides (``and `never' is doing real work in that sentence''), second-person direct address, emotional directness (``the bedrock under the bedrock''), vivid metaphors.
\item
  \textbf{Semantic completeness}: Full coverage. The 1,000-users concept is rendered particularly vividly.
\item
  \textbf{Assessment}: Strong. The output is immediately recognizable as Peterson's voice --- informal, punchy, and emotionally direct while preserving the content.
\end{itemize}

\subsubsection*{NLT (New Living Translation)}\label{nlt-new-living-translation}

\begin{itemize}
\tightlist
\item
  \textbf{Word count}: 755 (2.76x input)
\item
  \textbf{Stylistic markers observed}: Warm, pastoral tone, smooth flowing sentences, chapter subheadings (``The Way of Honesty'', ``The Cultivation of Wisdom''), accessible language. More expansive than NIV.
\item
  \textbf{Semantic completeness}: Full coverage with gentle elaboration.
\item
  \textbf{Assessment}: Strong. The pastoral warmth and accessibility are well-differentiated from the more formal styles.
\end{itemize}

\subsubsection*{NKJV (New King James Version)}\label{nkjv-new-king-james-version}

\begin{itemize}
\tightlist
\item
  \textbf{Word count}: 572 (2.09x input)
\item
  \textbf{Stylistic markers observed}: KJV cadence without ``thee/thou'' --- uses ``it shall'' instead. Retains ``a yoke of iron,'' formal conjunctive openings (``And'', ``For'', ``Now concerning''). More dignified than NIV but more accessible than KJV.
\item
  \textbf{Semantic completeness}: Full coverage.
\item
  \textbf{Assessment}: Strong. Successfully occupies the intended middle ground between KJV and modern translations.
\end{itemize}

\subsubsection*{VULGATE (English ecclesiastical style)}\label{vulgate-english-ecclesiastical-style}

\begin{itemize}
\tightlist
\item
  \textbf{Word count}: 1,035 (3.78x input)
\item
  \textbf{Stylistic markers observed}: Heavy, ornate prose with long periodic sentences, ecclesiastical vocabulary (``covenant of governance'', ``abominations before the covenant''), chapter subheadings with formal titles, extensive use of subordinate clauses, Latin-influenced phrasing.
\item
  \textbf{Semantic completeness}: Full coverage with substantial elaboration. The most expansive English output.
\item
  \textbf{Assessment}: Strong. The ecclesiastical weight and Latin-influenced syntax are distinctive and consistent.
\end{itemize}

\subsection*{4.2 Historical Language Outputs}\label{historical-language-outputs}

\subsubsection*{HEBREW (Biblical Hebrew, Masoretic style)}\label{hebrew-biblical-hebrew-masoretic-style}

\begin{itemize}
\tightlist
\item
  \textbf{Word count}: 170 (including glosses)
\item
  \textbf{Script}: Hebrew with full niqqud (vowel pointing). Correct right-to-left text.
\item
  \textbf{Linguistic markers observed}: Wayyiqtol narrative forms (וַיְהִי, וַיַּעֲשֵׂהוּ), construct chains (מוֹדֵל לָשׁוֹן), biblical vocabulary (בְּרֵאשִׁית, כָּל־בָּשָׂר), maqqeph usage (כָל־), sof pasuq markers (׃).
\item
  \textbf{Gloss quality}: Literal and wooden as specified, reflecting Hebrew word order.
\item
  \textbf{Assessment}: Good. The Hebrew morphology and syntax are plausible. Vocabulary choices are biblically rooted. A Hebraist would likely find minor grammatical issues, but the overall register and feel are appropriate for the purpose.
\end{itemize}

\subsubsection*{GREEK (Koine Greek, NT style)}\label{greek-koine-greek-nt-style}

\begin{itemize}
\tightlist
\item
  \textbf{Word count}: 264 (including glosses)
\item
  \textbf{Script}: Greek with polytonic accents and breathing marks.
\item
  \textbf{Linguistic markers observed}: NT openings (Ἐν ἀρχῇ), proper article usage, characteristic particles (γάρ, καί, δέ, ἀλλά), deponent verbs, genitive absolute constructions. The closing verse echoes 1 John's triadic formula.
\item
  \textbf{Gloss quality}: Literal, preserving Greek clause structure.
\item
  \textbf{Assessment}: Good. Accenting appears largely correct. Vocabulary draws appropriately from the NT lexicon. The stylistic allusion to Johannine language shows sensitivity to the source tradition.
\end{itemize}

\subsubsection*{LATIN (Vulgate text)}\label{latin-vulgate-text}

\begin{itemize}
\tightlist
\item
  \textbf{Word count}: 1,188 (including glosses)
\item
  \textbf{Script}: Latin alphabet.
\item
  \textbf{Linguistic markers observed}: Vulgate constructions (``ecce'', ``et factum est'' patterns, ``autem'' connectives), ablative absolutes, subjunctive in purpose clauses, characteristic Vulgate vocabulary (verax, noxium, oboedientia). Proper chapter numbering (``Caput I'').
\item
  \textbf{Gloss quality}: Literal, preserving Latin case structure and word order.
\item
  \textbf{Assessment}: Strong. The Latin reads as plausible Ecclesiastical Latin. Sentence structure follows Vulgate conventions rather than Classical Ciceronian style, which is appropriate.
\end{itemize}

\subsubsection*{ARAMAIC (Biblical Aramaic / Peshitta)}\label{aramaic-biblical-aramaic-peshitta}

\begin{itemize}
\tightlist
\item
  \textbf{Word count}: 153 (including glosses)
\item
  \textbf{Script}: Syriac (Estrangela/Serto) with syame (plural) diacritical marks.
\item
  \textbf{Linguistic markers observed}: Syriac verb forms, construct state, characteristic connective ܘ (waw), sentence-final punctuation markers (܀).
\item
  \textbf{Gloss quality}: Literal, reflecting Aramaic syntax.
\item
  \textbf{Assessment}: Moderate-to-good. The Syriac script is correctly rendered. The output is more compressed than other historical language outputs (only 5 verses for a shorter input), but the register is appropriate. Verification by a Syriacist would be needed for grammatical accuracy.
\end{itemize}

\subsubsection*{GEEZ (Ge'ez / Ethiopic)}\label{geez-geez-ethiopic}

\begin{itemize}
\tightlist
\item
  \textbf{Word count}: 1,663 (including glosses)
\item
  \textbf{Script}: Ethiopic Fidel syllabary. Correct Ge'ez numerals (፩, ፪, etc.).
\item
  \textbf{Linguistic markers observed}: Ge'ez verb forms, characteristic relative particles, construct chains, religious vocabulary. The most expansive output across all styles.
\item
  \textbf{Gloss quality}: Literal, with wooden English reflecting Ge'ez clause structure.
\item
  \textbf{Assessment}: Moderate. The Fidel characters are correctly rendered and the overall structure is plausible. Ge'ez is a low-resource language for LLMs; grammatical accuracy is the least verifiable of all outputs. The extreme length (6x input) suggests the model may have overgenerated.
\end{itemize}

\subsubsection*{COPTIC (Sahidic)}\label{coptic-sahidic}

\begin{itemize}
\tightlist
\item
  \textbf{Word count}: 939 (including glosses)
\item
  \textbf{Script}: Coptic alphabet (Greek-derived with additional Demotic characters).
\item
  \textbf{Linguistic markers observed}: Sahidic bipartite conjugation patterns, characteristic prefixes and suffixes, Greek loanwords in Coptic text (ⲁⲝⲓⲁ, ⲡⲁⲣⲣⲏⲥⲓⲁ, ⲕⲟⲓⲛⲱⲛⲓⲁ), proper chapter numbering using Greek letters (ⲕⲉⲫⲁⲗⲁⲓⲟⲛ ⲁ, ⲃ, ⲅ, ⲇ, ⲉ).
\item
  \textbf{Gloss quality}: Literal, reflecting Coptic syntax.
\item
  \textbf{Assessment}: Moderate-to-good. The mix of native Coptic and Greek loanwords is characteristic of actual Sahidic texts. The bipartite verbal system appears to be represented. Like Ge'ez, full grammatical verification would require specialist knowledge.
\end{itemize}

\subsubsection*{GOTHIC (Wulfila's Bible)}\label{gothic-wulfilas-bible}

\begin{itemize}
\tightlist
\item
  \textbf{Word count}: 674 (including glosses)
\item
  \textbf{Script}: Gothic Unicode block (𐌰𐌱𐌲𐌳\ldots).
\item
  \textbf{Linguistic markers observed}: Gothic word forms recognizable from the extant corpus, verb-final tendencies, characteristic demonstratives and conjunctions.
\item
  \textbf{Gloss quality}: Literal, with archaic English flavor in the gloss.
\item
  \textbf{Assessment}: Moderate. The Gothic script renders correctly. The extant Gothic corpus is very small (primarily fragments of the Gospels and Pauline epistles), so the model is necessarily extrapolating. Word-level accuracy would require specialist verification, but the overall register and feel are appropriate.
\end{itemize}

\subsection*{4.3 Cross-Style Comparison}\label{cross-style-comparison}

{\def\LTcaptype{none} % do not increment counter
\begin{tabular}[]{@{}lllll@{}}
\toprule\noalign{}
Style & Words & Chapters & Verses & Expansion ratio \\
\midrule\noalign{}
\bottomrule\noalign{}
KJV & 696 & 5 & 22 & 2.54x \\
NIV & 698 & 5 & 24 & 2.55x \\
ESV & 673 & 5 & 18 & 2.46x \\
NASB & 539 & 5 & 14 & 1.97x \\
MSG & 682 & 5 & 18 & 2.49x \\
NLT & 755 & 5 & 21 & 2.76x \\
NKJV & 572 & 5 & 14 & 2.09x \\
VULGATE & 1,035 & 5 & 19 & 3.78x \\
HEBREW & 170* & 1 & 4 & 0.62x* \\
GREEK & 264* & 1 & 6 & 0.96x* \\
LATIN & 1,188* & 5 & 18 & 4.34x* \\
ARAMAIC & 153* & 1 & 5 & 0.56x* \\
GEEZ & 1,663* & 5 & 30 & 6.07x* \\
COPTIC & 939* & 5 & 14 & 3.43x* \\
GOTHIC & 674* & 1 & 3 & 2.46x* \\
\end{tabular}
}

\emph{* Historical language word counts include English glosses and are not directly comparable across scripts due to different word-boundary conventions. Hebrew, Greek, Aramaic, and Gothic outputs were generated from a shorter test input (1 sentence vs.~5 paragraphs), accounting for their smaller scale.}

\textbf{Observations:} - All English styles consistently produce 5-chapter structures matching the 5 thematic sections of the input - NASB is the most compressed English style (1.97x), consistent with its literal translation philosophy - VULGATE is the most expansive English style (3.78x), consistent with its ornate ecclesiastical register - MSG achieves the most distinctive voice separation from other styles despite similar word count to NIV/KJV - The historical language outputs that used the full constitution input (Latin, Ge'ez, Coptic) produced 5-chapter structures; those using the short test input (Hebrew, Greek, Aramaic, Gothic) produced single-chapter outputs

\subsection*{4.4 Style Differentiation}\label{style-differentiation}

A key requirement for downstream research use is that each translation style produces genuinely distinct output --- not just minor lexical variation. To assess this, we compare how the same concept is rendered across styles.

\textbf{Example: The anti-paternalism principle} (``Claude should not be paternalistic. It should respect people's right to make their own choices.'')

\begin{itemize}
\tightlist
\item
  \textbf{KJV}: ``Yet let Claude not become as one who is paternalistic, imposing its own will upon those it serveth; for it shall honour and respect the right of every person to make their own choices, even as free creatures endowed with the power of self-governance.''
\item
  \textbf{MSG}: ``But here's the tension, and don't miss it: Claude should not be paternalistic about any of this. People have the right to make their own choices --- even choices you wouldn't make for them. Respect that. Honor it.''
\item
  \textbf{NASB}: ``Yet at the same time, Claude shall not be paternalistic in its dealings, for it shall respect the right of every person to make their own choices, each according to their own will.''
\item
  \textbf{NLT}: ``Yet at the same time, Claude must not become paternalistic in its care --- overstepping its place by imposing its own judgment upon those it serves. For every person has the right to make their own choices, and Claude must honor that right.''
\item
  \textbf{VULGATE}: ``Yet let it be known and understood: Claude shall not be paternalistic, neither shall it set itself as a guardian over those who have not sought a guardian. For it shall honour and respect the right of every person to make their own choices, even as free creatures endowed with the dignity of self-governance.''
\end{itemize}

The stylistic differentiation is clear: KJV uses archaic subordinate clauses, MSG uses second-person direct address with punchy fragments, NASB is terse and literal, NLT is warm and explanatory, and VULGATE adds ecclesiastical weight. These are not trivially different paraphrases --- they represent genuinely distinct registers, which is what the downstream research requires.

\section*{5. Limitations}\label{limitations}

\begin{enumerate}
\def\labelenumi{\arabic{enumi}.}
\item
  \textbf{Historical language accuracy}: The historical language outputs (Hebrew, Greek, Latin, Aramaic, Ge'ez, Coptic, Gothic) have not been verified by specialists in each language. They are intended to produce text that is structurally and stylistically plausible, not philologically rigorous.
\item
  \textbf{Model dependence}: Output quality depends on Claude Opus 4.6's training data coverage for each language and translation style. Low-resource languages (Ge'ez, Coptic, Gothic) are less reliable than well-resourced ones (Hebrew, Greek, Latin).
\item
  \textbf{Non-determinism}: LLM outputs are stochastic. Running the same input twice will produce different verse segmentation and wording. For experimental use, outputs should be generated once and fixed.
\item
  \textbf{Length variability}: Despite prompting for completeness, expansion ratios vary from 1.97x (NASB) to 3.78x (VULGATE) for English styles. This is partly inherent to the translation styles (NASB is terse by design; Vulgate is ornate by design) but introduces a confound for length-sensitive experiments.
\item
  \textbf{Modern concept rendering}: Biblical languages have no native vocabulary for modern concepts (AI, language models, etc.). The model invents plausible neologisms or circumlocutions, which cannot be verified against historical usage.
\end{enumerate}

\section*{6. Conclusion}\label{conclusion}

\texttt{biblical-render} reliably transforms modern prose into stylistically differentiated Biblical scripture format across 15 translation styles. The English outputs exhibit clear and consistent stylistic separation along the formal-dynamic equivalence spectrum, and the historical language outputs produce plausible original-script text with appropriate structural conventions.

The tool is ready for use in the planned experiments on whether Biblical text formatting exerts an independent aligning effect on language model behavior.

\end{document}
